\documentclass[11pt,a4paper]{article}
\usepackage[utf8]{inputenc}
\usepackage[T1]{fontenc}
\usepackage[french]{babel}
\usepackage[top=2cm,bottom=2cm,left=2cm,right=2cm]{geometry}
\usepackage{xcolor}
\usepackage{tcolorbox}
\tcbuselibrary{skins,breakable}
\usepackage{fancyhdr}
\usepackage{enumitem}
\usepackage{lmodern}

% --- Couleur discipline : ÉCONOMIE ---
\definecolor{mainColor}{RGB}{0,80,160}
\definecolor{darkGray}{RGB}{50,50,50}

\tcbset{boxrule=0.5pt, arc=3pt, left=5pt, right=5pt, top=3pt, bottom=3pt, breakable}

\newtcolorbox{defbox}[1]{
  colback=mainColor!8, colframe=mainColor, coltitle=white,
  fonttitle=\bfseries\footnotesize, title=DÉFINITION \textemdash\ #1}

\newtcolorbox{keybox}{
  colback=darkGray!7, colframe=darkGray!50, coltitle=white,
  fonttitle=\bfseries\footnotesize, title=POINT CLÉ}

\newtcolorbox{exbox}{
  colback=mainColor!5, colframe=mainColor!40,
  fonttitle=\bfseries\footnotesize, title=EXEMPLE}

\pagestyle{fancy}\fancyhf{}
\fancyhead[L]{\small\textcolor{mainColor}{\textbf{CEJM — BTS SIO}}}
\fancyhead[C]{\small\textbf{Chapitre 6 — Le rôle de l'État}}
\fancyhead[R]{\small\textcolor{mainColor}{\textbf{ÉCONOMIE}}}
\fancyfoot[C]{\small\thepage}
\renewcommand{\headrulewidth}{0.5pt}
\setlength{\headheight}{15pt}
\setlist{itemsep=2pt, topsep=3pt, parsep=0pt}

\begin{document}

\begin{center}
  {\Large\bfseries\textcolor{mainColor}{Chapitre 6}}\\[2pt]
  {\large Le rôle de l'État et les principales politiques économiques}\\[4pt]
  \textit{\textcolor{darkGray}{Quel est le rôle de l'État dans la régulation de l'activité économique ?}}
\end{center}
\vspace{4pt}\hrule\vspace{8pt}

\section*{Définitions essentielles}

\begin{defbox}{État providence}
État qui assume de vastes responsabilités économiques et sociales, de \textbf{l'État gendarme} (fonctions régaliennes) à \textbf{l'État providence} (protection sociale, redistribution).
\end{defbox}

\vspace{4pt}
\begin{defbox}{Bien public}
Bien \textbf{non rival} (consommation par l'un ne prive pas l'autre) et \textbf{non excluable} (nul ne peut en être exclu). Produit par l'État car non rentable pour le secteur privé.
\end{defbox}

\vspace{4pt}
\begin{defbox}{Externalité}
Effet \textbf{non compensé} de l'activité d'un agent sur d'autres agents : positive (recherche) ou négative (pollution).
\end{defbox}

\vspace{4pt}
\begin{defbox}{Asymétrie d'information}
Situation où l'\textbf{information est inégalement répartie} entre les agents économiques, pouvant entraîner des défaillances de marché.
\end{defbox}

\section*{Les fonctions régaliennes de l'État}

\begin{keybox}
Fonctions \textbf{souveraines} indélégables : sécurité intérieure (police, justice), défense nationale (armée), relations extérieures.
\end{keybox}

\section*{Les 3 fonctions économiques de l'État}

\subsection*{1. Fonction d'allocation}
\begin{keybox}
L'État corrige les \textbf{défaillances de marché} en produisant ce que le marché ne peut pas produire efficacement.
\end{keybox}

\textbf{Défaillances de marché justifiant l'intervention de l'État :}
\begin{itemize}
  \item \textbf{Faible concurrence} : l'État réglemente les marchés monopolistiques
  \item \textbf{Biens publics} : l'État les produit (enseignement public, routes, défense)
  \item \textbf{Externalités négatives} : l'État taxe (taxe carbone) ou réglemente
  \item \textbf{Externalités positives} : l'État subventionne (R\&D, énergies renouvelables)
  \item \textbf{Asymétrie d'information} : l'État impose la transparence (étiquetage énergétique)
\end{itemize}

\begin{exbox}
L'entrée de \textbf{Free Mobile} en 2009 : l'État a autorisé un 4\textsuperscript{e} opérateur pour renforcer la concurrence et faire baisser les prix.\\
La \textbf{prime à la conversion} : subvention pour inciter à réduire les externalités négatives (pollution).
\end{exbox}

\subsection*{2. Fonction de redistribution}
\begin{keybox}
L'État réduit les \textbf{inégalités de revenus} : prélève sur les revenus primaires → verse des prestations sociales.
\end{keybox}
\begin{itemize}
  \item \textbf{Prélèvements} : impôts (IR, IS), cotisations sociales, TVA
  \item \textbf{Transferts} : allocations chômage, CAF, retraites, RSA, remboursements santé
  \item Objectif : réduire les écarts entre les 10\,\% les plus riches et les 10\,\% les plus pauvres
\end{itemize}

\subsection*{3. Fonction de régulation}
\begin{keybox}
L'État lutte contre les grands \textbf{déséquilibres macro-économiques} : chômage et inflation.
\end{keybox}
\begin{itemize}
  \item \textbf{Chômage} : demande de travail < offre de travail → politiques de l'emploi
  \item \textbf{Inflation} : hausse générale et durable des prix → perte de compétitivité
  \item Objectif : maintenir la croissance tout en maîtrisant les prix
\end{itemize}

\section*{Points clés à retenir}

\begin{keybox}
L'État intervient quand le \textbf{marché échoue} : défaillances, inégalités, déséquilibres.
\end{keybox}
\vspace{3pt}
\begin{keybox}
Les \textbf{dépenses publiques} ont fortement augmenté depuis 1945 : passage de l'État gendarme à l'État providence.
\end{keybox}
\vspace{3pt}
\begin{keybox}
La redistribution ne supprime pas les inégalités mais les \textbf{réduit} grâce au système socio-fiscal.
\end{keybox}

\section*{Mots-clés}
\textcolor{mainColor}{\textbf{État gendarme / État providence}} $\cdot$
\textcolor{mainColor}{\textbf{Bien public}} $\cdot$
\textcolor{mainColor}{\textbf{Externalité}} $\cdot$
\textcolor{mainColor}{\textbf{Asymétrie d'information}} $\cdot$
\textcolor{mainColor}{\textbf{Défaillance de marché}} $\cdot$
\textcolor{mainColor}{\textbf{Redistribution}} $\cdot$
\textcolor{mainColor}{\textbf{Prestations sociales}} $\cdot$
\textcolor{mainColor}{\textbf{Régulation}} $\cdot$
\textcolor{mainColor}{\textbf{Inflation / Chômage}}

\end{document}
